% VI45-DETAILED-HINTS
% VI45-DETAILED-HINT-01
\paragraph{Exercise VI.45.1.}
\begin{hint}
% PEDAGOGY-ENRICHED-VI
Use the charts \(s\neq0\) and \(t\neq0\). Use the Rees-algebra or affine-chart description of the blow-up and track the exceptional divisor through the chart coordinates.
\end{hint}

% VI45-DETAILED-HINT-02
\paragraph{Exercise VI.45.2.}
\begin{hint}
% PEDAGOGY-ENRICHED-VI
Set \(x=y=0\). Use the Rees-algebra or affine-chart description of the blow-up and track the exceptional divisor through the chart coordinates.
\end{hint}

% VI45-DETAILED-HINT-03
\paragraph{Exercise VI.45.3.}
\begin{hint}
% PEDAGOGY-ENRICHED-VI
Substitute \(y=xu\) and remove the exceptional factor. Use the Rees-algebra or affine-chart description of the blow-up and track the exceptional divisor through the chart coordinates.
\end{hint}

% VI45-DETAILED-HINT-04
\paragraph{Exercise VI.45.4.}
\begin{hint}
% PEDAGOGY-ENRICHED-VI
Track the distinct points of \(E\) representing the two directions. Use the Rees-algebra or affine-chart description of the blow-up and track the exceptional divisor through the chart coordinates.
\end{hint}

% VI45-DETAILED-HINT-05
\paragraph{Exercise VI.45.5.}
\begin{hint}
% PEDAGOGY-ENRICHED-VI
Substitute \(y=xu\) and factor the maximal power of \(x\). Use the Rees-algebra or affine-chart description of the blow-up and track the exceptional divisor through the chart coordinates.
\end{hint}

% VI45-DETAILED-HINT-06
\paragraph{Exercise VI.45.6.}
\begin{hint}
% PEDAGOGY-ENRICHED-VI
Use the tangent-cone polynomial on \(E\simeq\mathbb P^1\). Use the Rees-algebra or affine-chart description of the blow-up and track the exceptional divisor through the chart coordinates.
\end{hint}

% VI45-DETAILED-HINT-07
\paragraph{Exercise VI.45.7.}
\begin{hint}
% PEDAGOGY-ENRICHED-VI
Use the relative-Proj definition. Use the Rees-algebra or affine-chart description of the blow-up and track the exceptional divisor through the chart coordinates.
\end{hint}

% VI45-DETAILED-HINT-08
\paragraph{Exercise VI.45.8.}
\begin{hint}
% PEDAGOGY-ENRICHED-VI
Apply the universal property to the identity. Use the Rees-algebra or affine-chart description of the blow-up and track the exceptional divisor through the chart coordinates.
\end{hint}

% VI45-DETAILED-HINT-09
\paragraph{Exercise VI.45.9.}
\begin{hint}
% PEDAGOGY-ENRICHED-VI
Use \(dH-mE\). Use the Rees-algebra or affine-chart description of the blow-up and track the exceptional divisor through the chart coordinates.
\end{hint}

% VI45-DETAILED-HINT-10
\paragraph{Exercise VI.45.10.}
\begin{hint}
% PEDAGOGY-ENRICHED-VI
Use the class \(H-E\). Use the Rees-algebra or affine-chart description of the blow-up and track the exceptional divisor through the chart coordinates.
\end{hint}

% VI45-DETAILED-HINT-11
\paragraph{Exercise VI.45.11.}
\begin{hint}
% PEDAGOGY-ENRICHED-VI
Use \(K_S=-3H+E\). Use the Rees-algebra or affine-chart description of the blow-up and track the exceptional divisor through the chart coordinates.
\end{hint}

% VI45-DETAILED-HINT-12
\paragraph{Exercise VI.45.12.}
\begin{hint}
% PEDAGOGY-ENRICHED-VI
Each point blow-up lowers \(K^2\) by one. Use the Rees-algebra or affine-chart description of the blow-up and track the exceptional divisor through the chart coordinates.
\end{hint}

% VI45-DETAILED-HINT-13
\paragraph{Exercise VI.45.13.}
\begin{hint}
% PEDAGOGY-ENRICHED-VI
Use the diagonal intersection form. Use the Rees-algebra or affine-chart description of the blow-up and track the exceptional divisor through the chart coordinates.
\end{hint}

% VI45-DETAILED-HINT-14
\paragraph{Exercise VI.45.14.}
\begin{hint}
% PEDAGOGY-ENRICHED-VI
Square \(2H-E_1-E_2-E_3\). Use the Rees-algebra or affine-chart description of the blow-up and track the exceptional divisor through the chart coordinates.
\end{hint}

% VI45-DETAILED-HINT-15
\paragraph{Exercise VI.45.15.}
\begin{hint}
% PEDAGOGY-ENRICHED-VI
Intersect with \(E_1\). Use the Rees-algebra or affine-chart description of the blow-up and track the exceptional divisor through the chart coordinates.
\end{hint}

% VI45-DETAILED-HINT-16
\paragraph{Exercise VI.45.16.}
\begin{hint}
% PEDAGOGY-ENRICHED-VI
Intersect with \(H-E_2-E_3\). Use the Rees-algebra or affine-chart description of the blow-up and track the exceptional divisor through the chart coordinates.
\end{hint}

% VI45-DETAILED-HINT-17
\paragraph{Exercise VI.45.17.}
\begin{hint}
% PEDAGOGY-ENRICHED-VI
Apply the Picard action twice. Use the Rees-algebra or affine-chart description of the blow-up and track the exceptional divisor through the chart coordinates.
\end{hint}

% VI45-DETAILED-HINT-18
\paragraph{Exercise VI.45.18.}
\begin{hint}
% PEDAGOGY-ENRICHED-VI
Use the full transformation for \((3;1,1,1)\). Use the Rees-algebra or affine-chart description of the blow-up and track the exceptional divisor through the chart coordinates.
\end{hint}

% VI45-DETAILED-HINT-19
\paragraph{Exercise VI.45.19.}
\begin{hint}
% PEDAGOGY-ENRICHED-VI
Use all four formulas for \((4;2,1,1)\). Use the Rees-algebra or affine-chart description of the blow-up and track the exceptional divisor through the chart coordinates.
\end{hint}

% VI45-DETAILED-HINT-20
\paragraph{Exercise VI.45.20.}
\begin{hint}
% PEDAGOGY-ENRICHED-VI
A point of \(E\) stores a tangent direction. Use the Rees-algebra or affine-chart description of the blow-up and track the exceptional divisor through the chart coordinates.
\end{hint}

% VI45-DETAILED-HINT-21
\paragraph{Exercise VI.45.21.}
\begin{hint}
% PEDAGOGY-ENRICHED-VI
Each pair met only at a blown-up vertex. Use the Rees-algebra or affine-chart description of the blow-up and track the exceptional divisor through the chart coordinates.
\end{hint}

% VI45-DETAILED-HINT-22
\paragraph{Exercise VI.45.22.}
\begin{hint}
% PEDAGOGY-ENRICHED-VI
Square \(dH-mE\). Use the Rees-algebra or affine-chart description of the blow-up and track the exceptional divisor through the chart coordinates.
\end{hint}

% VI45-DETAILED-HINT-23
\paragraph{Exercise VI.45.23.}
\begin{hint}
% PEDAGOGY-ENRICHED-VI
A Cartier divisor has an invertible ideal. Use the Rees-algebra or affine-chart description of the blow-up and track the exceptional divisor through the chart coordinates.
\end{hint}

% VI45-DETAILED-HINT-24
\paragraph{Exercise VI.45.24.}
\begin{hint}
% PEDAGOGY-ENRICHED-VI
Use \([uv:v:u]=[ut:t:1]\) on the blow-up chart. Use the Rees-algebra or affine-chart description of the blow-up and track the exceptional divisor through the chart coordinates.
\end{hint}
